% Generated draft: requires mathematical and source review.
\documentclass[11pt]{article}
\usepackage[margin=25mm]{geometry}
\usepackage{fontspec}
\setmainfont{Latin Modern Roman}
\usepackage{amsmath,amssymb,mathtools,unicode-math}
\setmathfont{Latin Modern Math}
\usepackage{enumitem,needspace,xurl,hyperref}
\setlength{\parindent}{0pt}
\setlength{\parskip}{5pt}
\setlength{\emergencystretch}{3em}
\newcommand{\N}{\mathbb{N}}
\newcommand{\Z}{\mathbb{Z}}
\newcommand{\Q}{\mathbb{Q}}
\newcommand{\R}{\mathbb{R}}
\newcommand{\C}{\mathbb{C}}
\newcommand{\F}{\mathbb{F}}
\newcommand{\A}{\mathbb{A}}
\newcommand{\PP}{\mathbb{P}}
\newcommand{\E}{\mathbb{E}}
\newcommand{\eps}{\varepsilon}
\newcommand{\dd}{\,\mathrm{d}}
\begin{document}
\section*{UnsolvedMath ID: 30003096}
\textbf{Source problem code:} catalogue entry
\subsection{Definitions and mathematical statement}
A discrete variational system is described schematically by a field \(u\) on a discrete multi-time space \(T\) and a discrete Lagrangian \(d\)-form \(\mathcal L[u]\). Here the space \(T\), the field's target space, the integer \(d\), and the allowed dependence of \(\mathcal L\) on the field are not specified in the supplied record. A discrete \(d\)-form assigns an action contribution to each oriented \(d\)-cell, changing sign when the orientation is reversed. For a finite oriented \(d\)-dimensional multi-time manifold \(\Sigma\) made of such cells, its action is schematically
\[
S_\Sigma[u]=\sum_{c\in\Sigma^{(d)}}\mathcal L[u](c),
\]
where \(\Sigma^{(d)}\) denotes its oriented \(d\)-cells. Infinite manifolds require an additional convergence or local-variation convention, which the source does not specify.

The stated pluri-Lagrangian property requires the same field \(u\) to be critical for every admissible \(\Sigma\). In first-variation notation, this is
\[
\delta S_\Sigma[u](v)=0
\qquad
\text{for every admissible }\Sigma
\text{ and every admissible variation }v.
\]
The field structure needed to define this first variation, the admissible manifolds, and the boundary restrictions on variations must be supplied before this becomes a fully specified variational problem. Fixed-boundary variations are a possible convention, but are not explicitly prescribed by the record.

\textbf{Problem.} What is integrability of discrete variational systems? More specifically, can integrability be characterized by the pluri-Lagrangian property that a field on a multi-time space is critical for the action of a \(d\)-form on every \(d\)-dimensional multi-time manifold?

The record does not specify a class of discrete variational systems or an independent definition of integrability against which to prove an equivalence. It also does not specify which fields must satisfy the simultaneous variational property, or how the multi-time system and its Lagrangian form are to be associated with a given system. Thus the displayed condition describes the proposed field-level property; it does not assert that the existence of one such field characterizes integrability. A precise necessary-and-sufficient theorem remains an unresolved formulation target.
\subsection{Short English statement}
The question asks whether integrability of a discrete variational system can be characterized by one field making the Lagrangian action stationary on every admissible multi-time surface of dimension \(d\). The system class, meaning of integrability, admissible surfaces, boundary conditions, and system-level quantifiers remain unspecified.
\subsection{Sources}
\begin{itemize}
\item Supplied local UnsolvedMath registry record, permanent ID \(30003096\), source problem code \texttt{catalogue entry}, in the Oberwolfach Reports Open Problems collection. The project registry is \texttt{data/unsolvedmath/problems.json}, with supplementary existing records in \texttt{catalog/problems.json}. Title: \emph{Pluri-Lagrangian Integrability of Discrete Variational Systems}. The record gives proposed year 2016 and names no individual proposer.
\item Supplied report citation: \emph{Geometric Numerical Integration} (2016), Oberwolfach Reports, pp.\ 869--948, DOI \texttt{10.4171/owr/2016/18}. \url{https://doi.org/10.4171/owr/2016/18}. The supplied background attributes the problem to workshop open-problem discussions documented in Oberwolfach Reports; it gives no problem-specific page or speaker attribution.
\item Supplied literature source: Boll--Petrera--Suris, DOI \texttt{10.1098/rspa.2013.0550}. \url{https://doi.org/10.1098/rspa.2013.0550}. No title or theorem locator is supplied.
\item Literature assessment dated exactly \texttt{2026-08-21}: the supplied curation labels the problem \emph{partially solved} and states that pluri-Lagrangian systems provide a successful notion and solve important one-dimensional and special lattice cases, while classification and the relation to standard integrability remain incomplete. This is supplied dated curation derived from the OWR record, not an independently verified literature assessment. Its time-sensitive status and cited claims require independent verification.
\end{itemize}
\end{document}
